

\section{Dinitz Algorithm}

Dinitz (or Dinic) Algorithm can solve max flow in $O(n^2 m)$ time. Here we first list some definitions:

\begin{definition}[Saturated Edge and Blocking Flow]
	An edge $e$ is called \textbf{saturated} if $f(e) = c(e)$. $f$ is called a \textbf{blocking flow} if every $s$-$t$ path contains at least one saturated edge.
\end{definition}

We can then make use of it to design Dinitz Algorithm. The main idea is to find a blocking flow in the residual graph, and then update the flow by adding the blocking flow. We repeat this process until there is no augmenting path in the residual graph.

\begin{algorithm}
	$f\gets 0$\;
	\While{there is an augmenting path in $G_f$}{
		$G' = (V,E')$, with $E' \subseteq E(G_f)$ being the edges in shortest path (assuming all edges have length 1) from $s$ in $G_f$\;
		$f' \gets$ a blocking flow in $G'$\;
		Update $f\gets f+f'$\;
	}
	\Return{$f$}\;
	\caption{Dinitz Algorithm}\label{alg:dinitz}
\end{algorithm}

\begin{lemma}
	In every iteration of Dinitz Algorithm, the length of the shortest path from $s$ to $t$ in the residual graph increases by at least 1.
\end{lemma}
\begin{proof}
	We can build layers based on the distance from $s$ in the residual graph. Then $G'$ is the subgraph after removing all edges that don't go from a lower layer to a higher layer (also note that all edges can only connect two neighboring layers). Since $f'$ is a blocking flow, there is no path that traverses all layers exactly once (i.e. $L_0 \to L_1 \to \cdots \to L_d$) in $G_{f+f'}$.

	As any newly created edges only go backwards, all paths will be strictly longer.
\end{proof}

Then we give an algorithm for computing blocking flows.

\begin{algorithm}[H]
	\While{there is an $s$-$t$ path in $G'$}{
		Apply DFS to find an $s$-$t$ path $p$ in $G'$\;
		Delete edges during backtracking, as they can't connect to $t$\;
		Push flow along $p$ by the bottleneck edge, and delete all saturated edges\;
	}
	\Return the flow\;
	\caption{Blocking Flow Algorithm}\label{alg:blocking}
\end{algorithm}

\begin{lemma}
	The blocking flow algorithm runs in $O(nm)$ time. If $\forall e\in E, c(e) =1$, then it runs in $O(m)$ time.
\end{lemma}
\begin{proof}
	Every iteration, we delete at least one edge, so there are at most $O(m)$ iterations. Each iteration takes $O(n)$ time to find a path and update the flow.

	If all edges have capacity 1, then each edge can only be saturated once, so the total time is $O(m)$.
\end{proof}

\begin{theorem}
	The Dinitz Algorithm runs in $O(n^2 m)$ time. If $c(s,u)$ and $c(v,t)$ are arbitrary, but all other edges have capacity 1, then it runs in $O(\min\{m^{3/2}, n^{2/3} m\})$ time.
\end{theorem}
\begin{proof}
	\begin{itemize}
		\item	As the distance from $s$ to $t$ increases by at least 1 in each iteration, there are at most $n-1$ iterations. Each iteration takes $O(nm)$ time, so the total time is $O(n^2 m)$.

		\item For the $m^{3/2}$ bound, we see that after $\sqrt{m}$ iterations, the distance from $s$ to $t$ is at least $\sqrt{m}+1$ and we have at least $\sqrt{m}+2$ layers, if we consider $s$ and $t$ as layers $L_0$ and $L_{\sqrt{m}+1}$ respectively.

		      We can find a cut between neighboring layers, so there are at least $\sqrt{m}$ cuts, where any edge can cross at most 1 cut. As the total capacity of all edges is $m$, we can guarantee that there is a cut with capacity at most $\sqrt{m}$. This means that the min cut of this graph (equivalently, max flow) is at most $\sqrt{m}$, so the total time is $O(m^{3/2})$ as we only need to iterate $\sqrt{m}$ times more.

		\item For the $n^{2/3} m$ bound, we see that after $3n^{2/3}$ iterations, the distance from $s$ to $t$ is at least $3n^{2/3}+1$ and we have at least $3n^{2/3}+2$ layers.

		      Then we see the average size of layers have size smaller than $n^{1/3}/3$, so at most $1/3$ of the layers have size greater than $n^{1/3}$ according to Markov's Inequality.

		      From Pigeonhole Principle, there are at least a pair of two neighboring layers with size smaller than $n^{1/3}$, so the cut between them has capacity at most $n^{2/3}$. This means that the min cut of this graph is at most $n^{2/3}$, and the rest follows similarly to the previous case.
	\end{itemize}
\end{proof}
\begin{remark}
	The $O(m^{3/2})$ bound can hold even when the graph is not simple, i.e. there can be multiple (a.k.a. parallel) edges between two vertices, but the $O(n^{2/3} m)$ bound won't hold there.
\end{remark}

\begin{theorem}
	A max matching in a bipartite graph can be found in $O(m\sqrt{n})$ time.
\end{theorem}
\begin{proof}
	We give direction to the edges, so that all edges are directed from $L$ to $R$. Then we can add two vertices $s$ and $t$,	adding $s$ connecting to all vertices in $L$, and $t$ connected by all vertices in $L$, with all edges having capacity 1. Then we can apply Dinitz Algorithm to find the max flow.

	After $2\sqrt{n}$ iterations, we have at least $2\sqrt{n}+2$ layers.

	Note that most of the edges are pointing from an odd layer to an even layer: the edges pointing from even layers to odd ones are `newly created' after pushing some flows and generating edges in the opposite direction, so they are indicating edges that are already matched.

	As the number of matched edges is at most $n$, the total number of `reverse' edges is at most $n$ as well. The number of cuts that cut from even layers to odd layers are $\sqrt{n}$, so there is a cut with capacity at most $n/\sqrt{n} = \sqrt{n}$. This means that the max flow is at most $\sqrt{n}$, so we only need to iterate $\sqrt{n}$ times more, which gives us the total time of $O(m\sqrt{n})$.
\end{proof}
